\section{Limitations}
\label{sec:limitations}

\paragraph{Scope of systems and data.}
We study three TDMs, two open backbones and two trained classifiers; frontier generative APIs are out of scope. All
five datasets are English, public and human-labelled, and the TREC-6 and FPB test sets are small (500 and 970
examples), which widens their intervals. Public benchmarks may appear in the pretraining or decision-training data of
the models, which we can check only partly (\hyperref[app:runs]{Appendix~\ref*{app:runs}}); contamination would favour the zero-shot systems.

\paragraph{Task format.}
All results concern candidate-set decisions. On CLINC150 and HWU64 each example offers the 18--22 labels of a fixed label
group, not all 150 or 64 intents; the groups were chosen to prevent leakage of the gold label through the option list.
Full-label classification is harder for zero-shot systems, and other groupings could change the absolute numbers; we did
not test alternative groupings.

\paragraph{Shift model.}
We simulate prevalence shift by resampling test sets, so the label-shift assumption holds by construction. Real
deployments also change $p(x\mid y)$ (covariate or concept shift), under which prior correction can fail; our results
are therefore an optimistic case for prior correction. The confirmatory tests use one shift level ($\alpha=1$); the
stronger level is reported descriptively. A naturally or temporally shifted dataset, and alternative candidate
groupings, would test how far our conclusions transfer; both require new inference runs and are left to future work.

\paragraph{Design of the shift tests.}
The confirmatory PA1--PA3 estimand is conditional on one fixed set of 50 target priors, and its intervals reflect only
the finiteness of the test set. In hindsight this was a weakness of our design. Single priors move the Jev--DeBERTa
difference by several points on FPB and HateXplain (standard deviation 5.4 and 7.4~points across priors), so the frozen
set of 50 does not pin down the mean over the shift distribution: on FPB the confirmatory PA1 estimate is 1.7~points,
the average over 2{,}000 priors 0.4 (\hyperref[tab:papop]{Table~\ref*{tab:papop}}). The significant FPB results for PA1 and PA2 may therefore
reflect the particular prior set as much as a property of Jev. PA3 is less affected: across prior sets its estimate keeps
its sign and stays between 1.7 and 2.7~points on FPB. Future studies with a fixed protocol should target the mean over the shift
distribution, with the priors resampled inside the bootstrap or with many more priors, rather than a single prior set.

\paragraph{Models.}
Jev's version cannot be pinned (every call returned \texttt{jev-1.13.0}), and Jev is not deterministic: its probability
vectors differ between repeated calls on up to 81\% of examples, although its predictions rarely do. We run the two 9B
models in 8-bit precision to fit a 12\,GB GPU, whereas Nimble is distributed for bf16 inference; both members of the pair
share the precision, so the comparison is matched, but the size of the adapter effect in bf16 is untested and the
probabilities of 8-bit models can depend on the computation path. The vendors' training
data are not public, so we can show where decision training helps or hurts (for Nimble-9B on FPB, on the neutral class)
but not why.

\paragraph{Statistics.}
Bootstrap intervals hold quantities fitted on the development set fixed, and so understate total uncertainty; they are
not adjusted for multiplicity, whereas the reported significance is. Absence of a significant difference is not
evidence of equivalence beyond the width of the reported intervals.
